% 相关工作章节

\section{相关工作}

\subsection{传统信用卡欺诈检测方法}

信用卡欺诈检测的早期工作主要依赖基于规则的风控引擎和传统统计方法。逻辑回归（Logistic Regression）因其简单性和可解释性被广泛使用，但其线性假设限制了捕捉复杂欺诈模式的能力。决策树和随机森林（Random Forest）通过集成多棵决策树提升了泛化能力，并能处理非线性关系。Bhattacharyya等人\cite{bhattacharyya2011data}比较了支持向量机（SVM）、随机森林和逻辑回归在欺诈检测上的表现，发现随机森林在精确率-召回率权衡上表现最佳。

\subsection{梯度提升方法}

梯度提升树（Gradient Boosting Decision Trees）已成为表格数据上的标杆方法。XGBoost\cite{chen2016xgboost}通过引入正则化项和列采样防止过拟合，在多个Kaggle欺诈检测竞赛中占据主导地位。LightGBM和CatBoost进一步优化了训练效率和类别特征处理。然而，这些模型在处理极高维稀疏特征和复杂特征交互时仍存在局限，且缺乏对特征重要性的细粒度解释。

\subsection{深度表格模型}

近年来，研究者开始探索将深度学习应用于表格数据的有效方法。Borisov等人\cite{borisov2024deep}在其综述中系统性地梳理了表格深度学习的两大技术路线：基于Transformer的架构（如FT-Transformer、SAINT）和基于显式特征交互的架构（如DeepFM、DCN-V2），并指出表格数据对深度学习的独特挑战——异构特征、缺失值、不规则分布以及缺乏旋转不变性等归纳偏置。以下概述与本论文直接相关的代表性工作。

\textbf{FT-Transformer}\cite{gorishniy2021revisiting}将Transformer架构适配到表格数据。核心思想是将每个数值特征通过Feature Tokenizer映射为嵌入向量，在所有特征令牌前添加一个可学习的CLS令牌，然后通过标准Transformer编码器进行自注意力计算。CLS令牌的最终表示用于分类。FT-Transformer在多个表格数据基准上取得了有竞争力的结果，但在中等规模数据上树模型仍可能保持优势\cite{grinsztajn2022tree}。

\textbf{DeepFM}\cite{guo2017deepfm}结合了因子分解机（Factorization Machine）和深度神经网络。FM组件通过嵌入向量的内积高效捕获二阶特征交互，Deep组件通过多层感知机学习高阶特征表示。两者共享特征嵌入层，实现端到端联合训练。DeepFM最初用于点击率预测，但其特征交互建模能力同样适用于欺诈检测。

\textbf{SAINT}\cite{somepalli2022saint}在FT-Transformer的基础上引入了两种注意力机制的混合——对特征令牌间的自注意力（列注意力）和对样本间的行注意力——并通过自监督对比预训练提升了小样本场景下的性能。SAINT的结构设计表明，注意力机制的双重应用是表格深度学习的重要方向。

\subsection{后FT-Transformer时代：2022--2025年进展}

尽管FT-Transformer在多个基准上取得了优异表现，Grinsztajn等人\cite{grinsztajn2022tree}在45个表格数据集上的大规模实证研究发现了一个重要反例：在中等规模数据（$<$10万样本）上，基于树的模型（尤其是XGBoost和随机森林）即使不经过超参数调优，仍然持续优于深度模型。该研究指出，树模型的轴对齐分裂策略天然契合表格数据缺乏旋转不变性的特点，而神经网络对无信息特征（随机噪声列）更为敏感，性能下降幅度显著更大。这一发现对深度表格模型的研究方向产生了深远影响。

TabPFN\cite{hollmann2023tabpfn}提出了范式性的突破：通过在数百万合成数据集上预训练一个Transformer，使其学会贝叶斯推理过程，从而能在0.1秒内对新的小规模表格分类问题（$\le$1000样本、$\le$100特征）做出预测，且无需任何训练。但TabPFN目前受限于输入维度和小样本规模，尚无法直接应用于信用卡欺诈检测这类大规模场景。

Dastidar等人\cite{dastidar2024survey}在信用卡欺诈检测领域对机器学习和深度学习方法进行了系统综述，发现深度模型（尤其是LSTM和自编码器）在高维特征建模和序列行为检测方面具有独特优势，但在实际部署中仍面临数据隐私、类别不平衡和概念漂移等核心挑战。该综述进一步强调了AUPRC而非ROC-AUC作为主要评估指标的合理性，以及将可解释性分析纳入模型开发流程的必要性。

综上所述，本文的研究定位如下：在15万+样本的信用卡欺诈数据集上，以FT-Transformer（全局自注意力）和DeepFM（显式特征交互）代表两种不同的归纳偏置，系统比较深度表格模型与传统梯度提升方法在AUPRC和商业成本指标上的表现。在此基础上，本文引入Tree-Supervised知识蒸馏方案，将XGBoost的树结构先验迁移至深度模型，探索"以树为师"的归纳偏置增强路径。同时，通过SHAP分析和注意力可视化，为不平衡学习和可解释性提供实证证据。

\subsection{类别不平衡处理}

针对类别不平衡问题，现有方法可分为三大类：（1）数据层面方法，如随机过采样和SMOTE\cite{chawla2002smote}，通过生成少数类合成样本来平衡训练集；（2）算法层面方法，如类别加权和代价敏感学习，通过调整损失函数中不同类别的权重来关注少数类；（3）损失函数改进，如Focal Loss\cite{lin2017focal}，通过降低易分类样本的损失贡献，使模型聚焦于难分类样本。本文在实验部分对这三类方法进行了系统比较。

\subsection{模型可解释性}

在金融风控场景中，模型可解释性不仅是监管合规的要求，更是获取业务方信任的关键。SHAP\cite{lundberg2017unified}基于Shapley值的博弈论框架，为每个特征分配对预测结果的贡献值，是目前最被广泛接受的可解释性方法之一。对于树模型，TreeExplainer提供了快速精确的SHAP值计算。对于深度模型，FT-Transformer的内置注意力机制提供了额外的可解释性维度——通过可视化[CLS] Token对各特征的注意力权重，可以直观理解模型在做出预测时关注了哪些特征。
